\section{Exact quotient eigenspaces for prescribed cross graphs}\label{sec:frame}

This section proves Theorem~\ref{thm:frame} and
Corollary~\ref{cor:frame-recovery}. Its hypotheses and conclusion
differ from those of the main recovery theorem: all internal degrees
are at least three, all admissible block sizes are allowed, and cross graphs can be
arbitrary prescribed simple biregular graphs. The conclusion identifies
whole real eigenspaces, rather than only their sign vectors.

\subsection{Strong mixing directly in the configuration law}
Initially let each $C_i$ be an independent degree-$d_i=d_{ii}$ uniform
configuration adjacency matrix on $N$ vertices, with loops counted
twice. The only parity conditions are $d_iN$ even. Fix an odd prime
$p$, put $s(x)=\one_N^Tx$, and define
\begin{equation}\label{eq:q-reference}
 q_i(x,r)=\begin{cases}
 \one_{\{r=d_ic\one_N\}},&x=c\one_N,\\
 p^{-(N-1)}\one_{\{s(r)=d_is(x)\}},&x\text{ nonconstant}.
 \end{cases}
\end{equation}
Let $P_i(x,r)=\Prob(C_ix=r)$ and
\begin{equation}\label{eq:mix-error}
 e_i(x)=\sup_r|P_i(x,r)-q_i(x,r)|,\qquad
 \eps_N=\max_i\sum_xe_i(x).
\end{equation}
Theorem~\ref{thm:config-mixing}, proved in
Appendix~\ref{app:configuration-mixing}, gives $\eps_N\to0$ at
fixed $p$ and fixed frame along every admissible size sequence.
The sum of maximum-atom errors is unscaled. Neither $p\nmid d_i$
nor $p\nmid N$ is required. No matching-sum moment is transferred
through contiguity here; we shall condition this actual configuration
moment directly on simplicity.

\begin{lemma}[Replacement for coupled targets]\label{lem:replacement}
For arbitrary target functions $r_i(x_1,\ldots,x_k)$, replacing all
$P_i(x_i,r_i)$ by $q_i(x_i,r_i)$ changes their summed product by at most
\begin{equation}\label{eq:replacement}
 k\eps_N(2p+\eps_N)^{k-1}=o_{p,Q}(1).
\end{equation}
The bound is uniform over every choice of target functions.
\end{lemma}
\begin{proof}
The summed maximum atoms satisfy
\[
 \sum_x\sup_r q_i(x,r)
 =p+(p^N-p)p^{-(N-1)}\le2p,
 \qquad
 \sum_x\sup_r P_i(x,r)\le2p+\eps_N.
\]
Telescope the product. In a term replacing factor $i$, bound its
absolute difference by $e_i(x_i)$ and each other factor by its
maximum atom, which depends only on its own $x_j$. The unrestricted
sum then factorizes. Its bound is
$\eps_N(2p+\eps_N)^{k-1}$ per replacement.
No reference kernel is being treated as a random linear map.
\end{proof}

\subsection{The reference first moment}
For $\theta\in\Fp$ put
$\ell_p=\dim_{\Fp}\ker(Q-\theta I)$. Independence gives
\[
 \E|\ker_{\Fp}(A-\theta I)|
 =\sum_{x_1,\ldots,x_k}
    \prod_iP_i\!\left(x_i,\theta x_i-\sum_{j\ne i}B_{ij}x_j\right).
\]
Use Lemma~\ref{lem:replacement} and split the reference sum according
to which $x_i$ are constant.

If every block is nonconstant, the $i$th reference condition is
\[
 d_{ii}s(x_i)=\theta s(x_i)-\sum_{j\ne i}d_{ij}s(x_j),
 \quad\text{or equivalently}\quad (Q-\theta I)s=0.
\]
For each prescribed sum, there are exactly $p^{N-1}$ vectors, at
most $p$ of which are constant. This remains true when $p\mid N$.
Thus the all-nonconstant contribution is between
$p^{\ell_p}(1-p^{2-N})^k$ and $p^{\ell_p}$, and is
$p^{\ell_p}+o(1)$. If every block is constant, write
$x_i=c_i\one_N$; the equations are exactly $(Q-\theta I)c=0$.
This contribution is exactly $p^{\ell_p}$.

For mixed patterns we use a rank bound which holds in every field.
\begin{lemma}[Disjoint row supports]\label{lem:cross-rank}
If $B$ is a zero-one square matrix of order $N$ with all row and
column sums $d\ge1$, then over every field
\begin{equation}\label{eq:cross-rank}
 \rank B\ge\frac{N}{1+d(d-1)}.
\end{equation}
\end{lemma}
\begin{proof}
Each row support meets at most $d(d-1)$ other row supports. Greedily
select at least $N/[1+d(d-1)]$ rows with pairwise disjoint nonempty
supports. They are linearly independent over every field, since each
has a coordinate equal to one absent from the other selected rows.
\end{proof}

Let $S$ be the nonempty proper set of nonconstant blocks and $T$ its
complement. Connectedness supplies an edge $ij$ with $i\in T$,
$j\in S$, and $d_{ij}>0$. The deterministic equation at constant
block $i$ forces
\begin{equation}\label{eq:mixed-equation}
 \sum_{r\in S}B_{ir}x_r
 =\left((\theta-d_{ii})c_i-
           \sum_{r\in T\setminus\{i\}}d_{ir}c_r\right)\one_N.
\end{equation}
The concatenated coefficient matrix contains $B_{ij}$ and hence has
rank at least $\gamma N$, with
\[
 \gamma=\min_{i<j:d_{ij}>0}[1+d_{ij}(d_{ij}-1)]^{-1}>0.
\]
For fixed constants, \eqref{eq:mixed-equation} permits at most
$p^{|S|N-\gamma N}$ old-vector tuples. Each nonconstant reference
factor is at most $p^{-(N-1)}$ and each constant factor at most one.
The whole pattern contributes at most $p^{k-\gamma N}$, after summing
the $p^{|T|}$ constants. Ignore all other equations, which can only
reduce this count. There are $2^k-2$ mixed patterns; when $k=1$ there
are none.

Combining the three cases proves the sharp configuration-model moment
\begin{equation}\label{eq:frame-moment}
 \E|\ker_{\Fp}(A-\theta I)|
 =2p^{\ell_p}+o_{p,Q}(1),
\end{equation}
uniformly over all permitted cross matrices and over $\theta\in\Fp$.
More explicitly, for $k>1$ the absolute error is at most
\[
 kp^{\ell_p}p^{2-N}+(2^k-2)p^{k-\gamma N}
       +k\eps_N(2p+\eps_N)^{k-1}.
\]
This is a first moment in the configuration model. Below we condition
the nonnegative excess over the deterministic old kernel directly
on simplicity; a qualitative contiguity transfer is unnecessary.

\subsection{Reduction of every fixed real target}
For integer $\lambda$, put $\ell=\nullity_\Q(Q-\lambda I)$.
Outside finitely many primes, quotient rank is preserved: all larger
integer minors vanish, and one nonzero minor of order $k-\ell$
remains nonzero modulo $p$. The old modular kernel has dimension
$\ell$. Real nullity at least $\ell+1$ makes all minors of order
$kN-\ell$ of $A-\lambda I$ vanish, forcing modular nullity at least
$\ell+1$ for every prime. If $K_p$ is the modular nullity, the
deterministic lower bound $K_p\ge\ell$ gives
\begin{equation}\label{eq:subtracted-markov}
 \Prob(K_p\ge\ell+1)
 \le\frac{\E p^{K_p}-p^\ell}{p^{\ell+1}-p^\ell}.
\end{equation}
Equation~\eqref{eq:frame-moment} yields a limiting upper bound
$1/(p-1)$ on real failure, uniformly in the cross matrices.
First let $N\to\infty$ at fixed good odd $p$, then take arbitrarily
large such primes. This proves the integer-target assertion.

We give the elementary algebraic extension to avoid an unnecessary
integer-eigenvalue restriction.
\begin{lemma}[Root primes]\label{lem:root-primes}
For every nonconstant $f\in\Z[t]$ there are arbitrarily large primes
at which $f$ has a root, avoiding any prescribed finite set.
\end{lemma}
\begin{proof}
Choose an integer $c=f(b)\ne0$ and enlarge the forbidden set to
include two and the prime divisors of $c$. Let $M$ be their product.
Expansion gives $f(b+cMt)=c h(t)$, where $h\in\Z[t]$ is nonconstant
and is congruent to one at every forbidden prime. Choose an integer
$t$ with $|h(t)|>1$. A prime divisor of $h(t)$ lies outside that set
and makes $b+cMt$ a residue root of $f$. This is the elementary
Schur prime-divisor argument; it asserts no density or effective rate.
\end{proof}

\begin{lemma}[Good reductions of the quotient]\label{lem:good-roots}
Let $\lambda$ be a real algebraic integer, $f$ its monic irreducible
minimal polynomial, and $\ell=\nullity_\R(Q-\lambda I)$. There is a
nonzero integer $m$ such that for every prime $p\nmid m$ and every
root $\theta\in\Fp$ of $f$, one has
$\nullity_{\Fp}(Q-\theta I)=\ell$.
\end{lemma}
\begin{proof}
Each minor of $Q-tI$ has integer coefficients. Minors of order
$k-\ell+1$ vanish at $\lambda$, hence are divisible by $f$ in
$\Z[t]$. For $\ell=0$ that assertion is vacuous. Choose a minor
$g(t)$ of order $k-\ell$ nonzero at $\lambda$; use the empty minor
$g=1$ when $\ell=k$. Irreducibility gives $\gcd(f,g)=1$ over
$\Q[t]$. Clearing denominators in B\'ezout's identity yields
$Uf+Vg=m\ne0$ with $U,V\in\Z[t]$. At every residue root of $f$
outside the prime divisors of $m$, the chosen minor is nonzero and
all larger minors vanish, proving the claim.
\end{proof}

If $\nullity_\R(A-\lambda I)\ge\ell+1$, every minor of order
$kN-\ell$ of $A-tI$ vanishes at $\lambda$. Its polynomial is
divisible by the monic $f$ in $\Z[t]$, so at every residue root
$\theta$ it vanishes modulo $p$. Hence modular nullity is at least
$\ell+1$. Lemma~\ref{lem:good-roots} identifies the old modular
dimension at the good primes; \eqref{eq:subtracted-markov} applies
with \eqref{eq:frame-moment}. Lemma~\ref{lem:root-primes} supplies
arbitrarily large good odd root primes, completing the same order of
limits. Finally, a real number which is not an algebraic integer
cannot be an eigenvalue of any integer matrix, including $A$ or $Q$.
This proves the fixed-real-target assertion in the configuration model.

\subsection{Direct simple-model moments and uniform quantifiers}
Let $\mathcal S$ be simplicity of all independent internal factors.
Its probability satisfies
\[
 s_N\longrightarrow s_Q:=
 \exp\!\left[-\frac14\sum_i(d_{ii}^2-1)\right]>0.
\]
Conditioning gives exactly independent uniform simple internal
factors. For fixed $p,\theta$, write
$K=|\ker_{\Fp}(A-\theta I)|$ and
$\ell_p=\nullity_{\Fp}(Q-\theta I)$. The deterministic old kernel
has $p^{\ell_p}$ elements. Hence
\begin{equation}\label{eq:frame-simple-moment}
 \E_{\rm simple}K
 \le p^{\ell_p}+
       \frac{\E_{\rm config}K-p^{\ell_p}}{s_N},\qquad
 \limsup_N\sup_B\E_{\rm simple}K
 \le(1+s_Q^{-1})p^{\ell_p}.
\end{equation}
The last inequality follows from \eqref{eq:frame-moment} and is a
raw moment bound in the actual simple law. It does not assert that
the configuration coefficient two remains the simple-model limit.
For the good primes and residue roots in the real-target reduction,
subtract the old contribution in Markov as in
\eqref{eq:subtracted-markov}. The resulting limiting real failure
bound is $s_Q^{-1}/(p-1)$, which tends to zero as arbitrarily large
fixed good primes follow the size limit. This proves
Theorem~\ref{thm:frame} along every admissible sequence.

All configuration estimates were uniform in the prescribed cross
matrices and the conditioning probability is independent of them,
so their supremum is controlled directly. Averaging allows any joint
random cross law independent of the internal factors. This does not
mean simultaneous success for every cross choice on one realization,
or allow cross matrices chosen after observing the internals.

For the cubic pair $(2,1)$, Theorem~\ref{thm:cubic-raw} supplies
\eqref{eq:raw-simple-moment}. Real nullity at least two forces
modular nullity at least two for every prime; subtracting the planted
$p$ vectors gives limiting failure at most $e^{3/2}/(p-1)$.
Large fixed primes after the size limit give the planted real line
at every $N$. Gauge invariance makes uniformity over the prescribed
cross matching exact. Together with the $a\ge3$ case above, this
proves planted-line simplicity for every fixed $a>b>0$, every
admissible size, and every prescribed simple cross matrix independent
of the internal randomness. The argument gives no effective rate.

\subsection{Rational row clustering recovers the whole frame}
Let $\chi_Q$ be the characteristic polynomial and put
\begin{equation}\label{eq:mQ}
 m_Q(t)=\frac{\chi_Q(t)}{\gcd(\chi_Q(t),\chi'_Q(t))},
\end{equation}
with monic gcd. Because $Q$ is symmetric and integral, this is its
minimal polynomial, is squarefree, and belongs to $\Z[t]$.
Its roots are all the distinct real quotient eigenvalues and its
degree is at most $k$.
Always $m_Q(A)J=Jm_Q(Q)=0$. Apply Theorem~\ref{thm:frame} to the
finitely many roots of $m_Q$. Real diagonalizability of $A$ gives,
with probability $1-o(1)$ uniformly in the cross matrices,
\begin{equation}\label{eq:old-polynomial}
 \ker_\R m_Q(A)
 =\bigoplus_{\lambda\in\spec Q}\ker_\R(A-\lambda I)=W.
\end{equation}
Compute a rational basis matrix $V$ for this kernel. On this event
$V=JT$ for an invertible $k\times k$ matrix $T$.
Its rows are equal exactly within each planted class: distinct rows
of $T$ cannot coincide because $T$ is invertible. Grouping equal
rational rows recovers the classes without choosing individual
irrational eigenvectors or splitting repeated eigenspaces.

All coefficients and degrees of $m_Q$ are fixed. If the maximum
total degree of $A$ is $\Delta_Q$ and $m_Q(t)=\sum c_rt^r$, every
absolute row sum of $m_Q(A)$ is at most
$\sum_r|c_r|\Delta_Q^r=O_Q(1)$. Its minors are therefore bounded by
$\exp(O_Q(kN))$ by Hadamard's inequality. Forming this fixed-degree
matrix polynomial and using fraction-free rational elimination have
polynomial bit complexity. This proves
Corollary~\ref{cor:frame-recovery}.

Within the promised frame model, finding exact nullity
$\nullity_\Q m_Q(A)=k$ certifies \eqref{eq:old-polynomial}, because
the old $k$-dimensional space is always included. The algorithm assumes
the frame is known; unknown-frame identifiability is a different problem.

For the homogeneous frame $Q=(a-b)I_k+b\one_k\one_k^T$, with
$a\ge3$, $b\ge1$, and $k>1$, its two distinct eigenvalues are
$a+(k-1)b$ and $a-b$, the latter of multiplicity $k-1$ when $k>1$.
Then $m_Q(t)=(t-a-(k-1)b)(t-a+b)$ for $k>1$. This supplies recovery
for every fixed number of equal classes at all admissible sizes, without an assortativity
or outlier condition. In particular $k=2$ gives the $a\ge3$ part
of Corollary~\ref{cor:two-frame}, without an assortativity condition.
